% Section 14: The exterior derivative in geometric algebra} | Part: 2 | Depends on: see docs/REWRITE_PLAN.md | Status: todo | Original lines: 2289-2300
%!TEX root = ../../../main.tex
% =====================================================================
\section{The exterior derivative in geometric algebra}\label{sec:ga-calculus}
% =====================================================================

The exterior derivative now returns in a second language. The explicit coordinate formula \eqref{eq:dcoord} is the ``wedge half'' of one single first-order operator on multivector fields, the \emph{vector derivative}
$\vd=e^i_\sharp\,\partial_i$ (bold nabla, not to be confused with the covariant derivative $\nabla$). Differential forms are identified with multivector fields through $\flat$ (Proposition~\ref{prop:lcisiota}).
Let $V=\R^{p,q}$ with an orthonormal frame $(e_i)$ and reciprocal frame $(e^i)$, $e^i\cdot e_j=\delta^i_j$ (so $e^i=\varepsilon_ie_i$). On $V$ itself, coordinates $x^i$ are defined by $x=x^ie_i$. Using the metric, identify a differential form with a multivector field via $\dd x^i\leftrightarrow e^i$:
\[
 \omega=\sum_I\omega_I\,\dd x^I\ \longleftrightarrow\ F=\sum_I\omega_I\,e^{i_1}\wedge\cdots\wedge e^{i_k}.
\]
% Section 14: 01-vector-derivative | Part: 2 | Depends on: see docs/REWRITE_PLAN.md | Status: todo | Original lines: 2301-2335
%!TEX root = ../../../main.tex
\begin{definition}[Vector derivative]
The \emph{vector derivative} (Dirac operator) is $\vd=e^i\partial_i$. Its action on a multivector field $F$ is its geometric product with $F$:
\[
 \vd F=e^i\,\partial_iF=\vd\lc F+\vd\wedge F,\qquad \vd\lc F:=e^i\lc\partial_iF,\quad \vd\wedge F:=e^i\wedge\partial_iF .
\]
\end{definition}

For a $k$-vector field $F$, $\vd\wedge F$ has grade $k+1$ (``curl-like'') and $\vd\lc F$ has grade $k-1$ (``divergence-like''); this is formula \eqref{eq:aA} with the vector $a$ replaced by the operator $\vd$.

\begin{theorem}[$\dd$ is $\vd\wedge$]\label{thm:dGA}
Under the identification above:
\begin{enumerate}
  \item $\dd\omega\ \longleftrightarrow\ \vd\wedge F$.
  \item (Euclidean case) the codifferential $\codiff=(-1)^{n(k+1)+1}\hodge\dd\hodge$ corresponds to $-\vd\lc F$; hence $\vd=\dd-\codiff$ acting on forms.
  \item $\vd\wedge(\vd\wedge F)=0$ and $\vd\lc(\vd\lc F)=0$.
  \item $\vd^2=\vd\cdot\vd=\sum_i\varepsilon_i\partial_i^2$ (the Laplacian, or d'Alembertian), and $\vd^2=\vd\lc(\vd\wedge\,\cdot\,)+\vd\wedge(\vd\lc\,\cdot\,)=-(\dd\codiff+\codiff\dd)$ in the Euclidean case.
  \item (Duality) in terms of the dual of Definition~\ref{def:dual}:
  \[
   \dual{(\vd\wedge F)}=\vd\lc\dual F,\qquad \dual{(\vd\lc F)}=\vd\wedge\dual F .
  \]
\end{enumerate}
\end{theorem}
\begin{proof}
(1) The frame is constant, so $\partial_i$ commutes with it and
$\vd\wedge F=e^j\wedge\partial_jF=\sum_{I,j}\partial_j\omega_I\,e^j\wedge e^{i_1}\wedge\cdots\wedge e^{i_k}$, which is the image of \eqref{eq:dcoord}.
(2) The standard formula for the codifferential is $\codiff=-\sum_i\iota_{e_i}\partial_{i}$ in Euclidean space, and $\iota_{e_i}=e_i\lc$ on multivectors; so $\codiff=-\vd\lc$. Then $\vd=\vd\lc+\vd\wedge=-\codiff+\dd$.
(3) $\vd\wedge(\vd\wedge F)=e^i\wedge e^j\wedge\partial_i\partial_jF=0$ since $e^i\wedge e^j$ is antisymmetric and $\partial_i\partial_j$ symmetric (this is $\dd^2=0$ again). Similarly, by Proposition~\ref{prop:contr}(3), $\vd\lc(\vd\lc F)=(e^i\wedge e^j)\lc\partial_i\partial_jF=0$.
(4) $\vd\vd=e^ie^j\partial_i\partial_j=(e^i\cdot e^j)\partial_i\partial_j$ because the antisymmetric part $e^i\wedge e^j\partial_i\partial_j$ cancels. Expanding $\vd(\vd F)$ by grades and using (3) leaves the two cross terms, which equal $-(\dd\codiff+\codiff\dd)$ by (1)--(2).
(5) is Theorem~\ref{thm:dual}(3) applied with $A=\vd$ (the identities are linear in $A$ and the derivatives $\partial_i$ commute with the constant dual $\pss^{-1}$).
\end{proof}

\begin{remark}
Parts (2) and (4) say that the Dirac operator $\vd$ is a ``square root'' of the Laplacian, $\vd^2=-\Delta_{\mathrm{Hodge}}$, which splits as $\vd=\dd-\codiff$. Item (5) shows why the codifferential is $\pm\hodge\dd\hodge$: it is the \emph{dual} of the exterior derivative, just as contraction is the dual of the wedge.
\end{remark}

% Section 14: 02-examples | Part: 2 | Depends on: see docs/REWRITE_PLAN.md | Status: todo | Original lines: 2336-2365
%!TEX root = ../../../main.tex
\begin{example}[Vector calculus in $\GA(\R^3)$]
Let $f$ be a scalar field and $\mathbf a$ a vector field. Then
\[
 \vd f=\operatorname{grad}f,\qquad
 \vd\mathbf a=\vd\cdot\mathbf a+\vd\wedge\mathbf a=\operatorname{div}\mathbf a+\pss\,\operatorname{curl}\mathbf a ,
\]
using Proposition~\ref{prop:cross} ($\vd\wedge\mathbf a=\pss\,(\vd\times\mathbf a)$). Since $\operatorname{curl}\mathbf a=\dual{(\vd\wedge\mathbf a)}$ (Proposition~\ref{prop:cross}), the two classical identities are both instances of $\vd\wedge\vd\wedge{}=0$: $\operatorname{curl}\operatorname{grad}f=\dual{(\vd\wedge\vd f)}=0$, and, by \eqref{eq:dualduality},
$\operatorname{div}\operatorname{curl}\mathbf a=\vd\lc\dual{(\vd\wedge\mathbf a)}=\dual{(\vd\wedge\vd\wedge\mathbf a)}=0$. Finally $\vd^2=\vd\cdot\vd$ is the Laplacian.
\end{example}

\begin{example}[Maxwell's equations in one line]
In the spacetime algebra $\GA(\R^{1,3})$ the electromagnetic field is the bivector $F=\mathbf E+\pss\mathbf B$ (with $\mathbf E,\mathbf B$ the electric and magnetic fields written as relative vectors, i.e.\ as bivectors of spacetime) and (in suitable units) the four Maxwell equations are the single equation
\[
 \vd F=J ,\qquad\text{i.e.}\qquad \underbrace{\vd\wedge F=0}_{\dd F=0\ (\text{homogeneous})},\quad
 \underbrace{\vd\lc F=J}_{\text{inhomogeneous, source }J}.
\]
The two pieces are precisely the exterior derivative and the codifferential (a dual exterior derivative) of $F$. The first equation has a potential, $F=\vd\wedge A$, by the Poincar\'e lemma.
\end{example}

\begin{intuition}
In geometric algebra the derivative is \emph{literally a vector}, $\vd=e^i\partial_i$, so it multiplies a field $F$ with the geometric product, and the geometric product $ab=a\cdot b+a\wedge b$ splits into a lower-grade and a higher-grade part:
\[
 \underbrace{\vd F}_{\text{one derivative}}=\underbrace{\vd\lc F}_{\text{divergence-like (grade}\downarrow)}+\underbrace{\vd\wedge F}_{\text{curl-like = exterior derivative (grade}\uparrow)} .
\]
The exterior derivative is thus the ``wedge half'' of a single, invertible derivative. Duality then converts the wedge half into the contraction half, so that curl and divergence are \emph{the same operation seen through the dual} ($\hodge$): that is why Maxwell's equations, which look like two different kinds of equation in vector calculus, become one equation. Everything in this section is just $ab=a\cdot b+a\wedge b$ with $a$ replaced by $\vd$.
\end{intuition}

